\documentclass[11pt,a4paper]{article}
\usepackage[utf8]{inputenc}
\usepackage[T1]{fontenc}
\usepackage{lmodern}
\usepackage[margin=1in]{geometry}
\usepackage{graphicx}
\usepackage{hyperref}
\usepackage{booktabs}
\usepackage{longtable}
\usepackage{xcolor}
\usepackage{listings}
\usepackage{amsmath}
\usepackage{amssymb}
\usepackage{enumitem}
\usepackage{titling}
\usepackage{fancyhdr}
\usepackage{caption}
\usepackage[backend=biber,style=apa]{biblatex}

\addbibresource{references.bib}

\definecolor{moodgreen}{HTML}{2E7D5B}
\hypersetup{colorlinks=true, linkcolor=moodgreen, urlcolor=moodgreen, citecolor=moodgreen}

\pagestyle{fancy}
\fancyhf{}
\rhead{MoodBloom --- Group 2}
\lhead{\leftmark}
\rfoot{\thepage}

\title{CSC234 Final UI/UX Report --- MoodBloom}
\author{
  Kraiwich Jaiton \and
  Teerin Kittichaicharoen (UI/UX + QA Lead) \and
  Theerawat Patthawee (Project Lead) \and
  Jedsarit Fanpimiy \and
  Napat Chang-ekwong (UI/UX Lead) \\[1ex]
  \small Group 2 \\
  \small CSC234 User-Centric Mobile App Development \\
  \small KMUTT, Semester 2/2568
}
\date{May 28, 2026}

\begin{document}
\maketitle
\thispagestyle{empty}
\newpage

\noindent\textbf{Reviewer note.} Chapters~1 (User Research) and 2 (User Journey
Maps) are reconstructed from the persona patterns referenced in the Sprint 4--5
spec and the Sprint 5 kickoff prompt. Biographical details and direct quotes are
plausible reconstructions, not transcripts of interviews. The behaviours,
journey arcs, and acceptance-criteria mappings match the codebase's surface
decisions. \textbf{Requires team validation before submission.}

\tableofcontents
\newpage

\section*{Executive Summary}
\addcontentsline{toc}{section}{Executive Summary}

MoodBloom is a cross-platform Flutter mood tracker for Thai young adults that
uses AI to assist mood logging, detects distress patterns compassionately, and
visualises emotional history as a living garden. The UI/UX problem the team
set out to solve is the tension between gamification (which sustains
engagement) and clinical safety (which is non-negotiable when the population
includes users in genuine mental-health distress). The design philosophy that
resolved the tension is the \emph{ecosystem model}: every mood is weather,
plants are NEVER destroyed regardless of mood content, and therapeutic value
is never made contingent on mood state. Built across five sprints
(S1--S5) under a four-agent multi-agent workflow on Claude Code. The
release-candidate at submission passes 1018 automated tests including the
load-bearing TC-40 (Tier 3 must never call Gemini) and TC-41 (Safety
Filter rejects 100\% of off-script suggestions) with WCAG 2.2 AA contrast
across both themes \autocite{firth2017meta,nahum2018jitai}.

\section{User Research}

\subsection{Methodology}

The persona work was conducted as a stakeholder-grounded composite of two
user patterns the team identified during Sprint~1 agile planning. Persona
patterns were validated by mapping each acceptance test case in the Sprint
4--5 spec to a persona's expected behaviour. Direct user interviews were
not conducted within the academic timeline; the personas function as
design hypotheses anchored to published mental-health literature on Thai
young adults \autocite{firth2017meta,frattaroli2006expressive}.

\subsection{Persona 1: Lin --- harvest-oriented user}

22-year-old final-year university student. No diagnosed condition;
self-describes as ``thoughtful'' and uses journaling as a reflection tool.
Daily tech: mid-range Android, occasional web sessions. Goals: build a
sustainable mood-logging habit; notice weekly patterns; have something
gentle to look back at. Pain points: streak-shaming apps make her feel
worse on busy weeks; clinical-language apps feel diagnostic; ``fix-your-
mood'' apps feel performative. Acceptance-criteria coverage: TC-1..TC-5
(tokens), TC-11..TC-15 (harvest), TC-16..TC-20 (atmosphere),
TC-21..TC-24 (EWMA --- plants stay alive on bad weeks).

\subsection{Persona 2: Som --- user experiencing escalation}

27-year-old junior software engineer at a Bangkok startup. Two prior
work-stress-driven depressive episodes; one prior counsellor relationship.
Not on medication; comfortable with technology but wary of mental-health
apps that overpromise. Goals: notice heavy stretches early enough to act
on them; have a non-judgemental record; find a hotline quickly if needed.
Pain points: PHQ-9 questionnaires feel like diagnoses; aggressive crisis
surfacing trains ignorance; buried resources are useless when needed.
Acceptance-criteria coverage: TC-25..TC-30 (pattern detection),
TC-31..TC-35 (notification cooldown + Hotline 1323 in Tier 3 only),
TC-36..TC-39 (bipolar disclaimer), TC-40..TC-41 (Tier 3 determinism).

\section{User Journey Maps}

\subsection{Lin's evening logging + weekly harvest}

Five phases: Open ($<2$\,s cold start); Log (\texttt{AISuggestionPill} with
12-char minimum gate so 2--3-char drafts don't fire Gemini); Reflect
(missed days are empty slots, never streak breaks); Harvest (compassionate
copy ``a fresh canvas for your story''); Look back (24-hour mutability
window) \autocite{pennebaker1997expressive,locke2002goals}.

\subsection{Som's Tier 1 $\rightarrow$ 2 $\rightarrow$ 3 escalation}

Five phases: Drift (Mann--Kendall trend test fires; cooldown 1/24h, 48h
between alerts; Tier 1 routes to 2-minute breathing screen); Sustained low
(sliding 5-of-7; Tier 2 routes to journaling-prompt screen which reuses
\texttt{SaveMoodEntryUseCase}); Acute (3-consecutive $S \leq -0.6$ OR
$z_{\text{day}} < -2.5$ OR CUSUM breach; \textbf{Tier 3 messages come
from a curated pool only, never from Gemini} per ADR-0012); Opt-out
(advances cooldown anchor 48h so system does not re-nag); Hotline
(prominent tile, compassionate palette, never red alarm)
\autocite{stanley2012safety,nahum2018jitai}.

\section{Design Philosophy: The Ecosystem Model}

\subsection{The pivot}

In Sprints 1--3 the metaphor was ``flowers bloom for positive moods, wilt
for negative ones'' --- the standard wellness-app metaphor. After the
Sprint-3 demo the team's faculty advisor surfaced the clinical concern
that \emph{a user already in distress watching their plants wilt is
being shown a digital reinforcement of their state}. Sprint 4 replaced
``wilting plants'' with the ecosystem model:

\begin{quote}
Plants are NEVER destroyed, wilted, dying, or dead --- in any visual state,
in any copy, in any notification. Every mood is weather; the ecosystem
holds.
\end{quote}

Storm Season --- the worst of the five plant tiers --- paints plants as
\emph{sheltered}, with brighter lanterns and gentle rain falling around
them. The plants survive the storm.

\subsection{Therapeutic foundations}

\begin{description}
\item[Self-compassion] \autocite{neff2003selfcompassion,neff2023selfcompassion}: empirical case that self-kindness correlates with lower depression and anxiety. ``Missed days are empty slots, never streak breaks'' operationalises this.
\item[DBT validation] \autocite{linehan1993dbt}: emotions are valid information, not problems. Copy bans ``fix-your-mood'' verbs in favour of ``notice / explore / care for / pause.''
\item[ACT ``emotions as weather''] \autocite{hayes1999act}: emotions are transient phenomena. The metaphor's visual embodiment.
\item[Narrative externalisation] \autocite{white1990narrative,white2007maps}: separating person from problem. Weekly harvest copy frames the user as narrator.
\end{description}

\subsection{The three formulas}

Mood Score per entry, range $[-1, +1]$:
\[ S_t = v \cdot \frac{i}{5}, \qquad v \in \{-1, +1\}, \; i \in \{1, 2, 3, 4, 5\} \]

Daily Atmosphere drives sunny / calm / light-rain / storm:
\[ \overline{S}_{\text{today}} = \frac{1}{n} \sum_{i=1}^{n} S_i \]
Resets midnight. Storm visualises rain falling \emph{around} plants, never on them.

Garden Health EWMA smoothed across the week:
\[ H_t = \alpha \cdot S_t + (1 - \alpha) \cdot H_{t-1}, \qquad \alpha = 0.15, \; H_0 = 0 \]
The constant $\alpha = 0.15$ gives a $\sim$13-day window via
$\alpha \approx 2/(N+1)$, aligning with the PHQ-9's two-week period
\autocite{kroenke2001phq9,smit2022ewma}. The bound
$|\Delta H| \leq 0.15$ guarantees one bad day cannot crash the canvas.
Plant-tier mapping: $H \geq +0.4$ Flourishing; $+0.1 \leq H < +0.4$
Thriving; $-0.1 \leq H < +0.1$ Resting; $-0.4 \leq H < -0.1$ Weathering;
$H < -0.4$ Storm Season.

\section{Design System}

Tokens in \texttt{packages/design\_system/}. \texttt{MbColors} extension
extends Material 3's container palette with \texttt{text, textDim, bg,
card, line, skyBot, softCoral, softGreen}. Typography: Nunito for body,
Fraunces for display headings. Spacing tokens
(\texttt{MoodBloomSpacing.pagePadding, .radiusFull, .radiusSky}). Four
responsive breakpoints. Custom widgets: \texttt{MbCard, MbSectionLabel,
MbIntensityDots, MbSideNav}.

\section{Information Architecture \& Navigation}

GoRouter \texttt{StatefulShellRoute.indexedStack} for the four main tabs.
Modal routes (\texttt{/log-mood}, \texttt{/intervention/*},
\texttt{/privacy/setup}, \texttt{/unlock-history}) sit outside the shell.
Deep-link strategy: every \texttt{/history/:id} is gated by the History
privacy lock when the user has opted in (ADR-0013). The intervention
banner host wraps \texttt{MaterialApp.router}'s \texttt{builder:} so it
overlays any route.

\section{Key Screens}

Six surfaces drive the user experience: \emph{Onboarding} with the
bipolar disclaimer slide; \emph{Log Mood} with intensity slider + AI
suggestion pill; \emph{Garden} with the 5 tier states $\times$ 4
atmospheres (plants alive in every combination, with the tier pill +
``Take a breath'' + ``Patterns'' CTAs); \emph{Insights} with the
mandatory disclaimer ack on first view, mood-score chart with EWMA
dashed overlay, marker band, tier-band legend, recent-triggers list;
\emph{Settings} with the privacy section (biometric + PIN + WebAuthn
preview); \emph{History} with calendar + list view, gated behind the
privacy lock. Screenshots: pending --- see
\texttt{reports/images/README.md}.

\section{Gamification Ethics: Token Economy \& Skins}

Six anti-pattern guardrails per Cheng et al. (2019) and SDT
\autocite{cheng2019gamification,deci2000sdt,ryan2017sdt}:

\begin{enumerate}[leftmargin=*]
\item No premium currency, only one currency type.
\item No FOMO --- no expiring items, no limited-time pressure.
\item No leaderboards, no friend comparison.
\item No loss aversion --- tokens never decay.
\item Cosmetic only --- therapeutic features (analytics, pattern detection, interventions, disclaimer) are ALWAYS free. TC-35 verifies no \texttt{tokens/} import in \texttt{intervention/}.
\item Optional visibility --- user can hide the token UI.
\end{enumerate}

The \emph{mood-agnostic earning rule} is structurally enforced: a file-
level test greps \texttt{award\_daily\_tokens.dart} for any import of
\texttt{MoodType, MoodEntry, MoodScore} --- the file CANNOT see the user's
mood content. The contingent-rewards risk Cheng et al. (2019) flagged
cannot recur.

\section{Therapeutic Safety}

\textbf{Quote Safety Filter (TC-41):} Tier 1/2 use a Gemini hybrid;
suggestions pass through \texttt{QuoteSafetyFilterImpl.gate} which
enforces a 140-char cap, a forbidden-word blacklist (depression,
diagnose, must, should, $\ldots$), and a $\geq 80\%$ whitelist overlap.
55 adversarial inputs reject 100\%.

\textbf{Bipolar/medical disclaimer placement (TC-36--TC-39):} onboarding
slide; mandatory ack on first Insights view (one-way Firestore rule);
short footer on every Tier 1/2/3 notification body; Settings $\to$ About
contains the full text verbatim.

\textbf{Cooldown (TC-31, TC-32):} max 1 notification per 24h; 48h
cool-down between alerts; \texttt{InterventionStateRepository} persists
\texttt{lastTriggeredAt} to Firestore with a SharedPreferences mirror
(ADR-0008). Opt-out advances the cooldown anchor 48h so the system does
not re-nag.

The curated Tier 3 phrase pool contains eight entries, each including
``Hotline 1323'' (sample: ``We care about you. If it helps to talk, the
Thai Mental Health Hotline is free at 1323, 24 hours.''). The pool was
read aloud by the entire team before merge, per ADR-0012 §5.

\section{Accessibility}

WCAG 2.2 AA verified across both themes. 16 token pairs $\times$ 2
themes = 32 measurements. Highlights: body text on card 14.68:1 light /
12.08:1 dark (AAA); Tier 3 crisis-tile 7.20:1 (AAA both); Tier 1/2 banner
11.36:1 / 11.10:1 (AAA). One known LOW (textDim/softCoral at 4.38:1 light
pre-fix) fixed in commit \texttt{b864e438}. Dynamic-type at 200\% checked
on every dialog; Insights ack dialog uses \texttt{scrollable: true}.
Semantics labels on every interactive widget; breathing-screen animated
circle excluded from semantics with the \texttt{'Breathing rhythm
guide'} label.

\section{Cross-Platform Considerations}

\textbf{Android.} FCM; \texttt{local\_auth} biometric with
Keystore-backed credential persistence; biometric gate route guard
(ADR-0008 + ADR-0013). Min-SDK 21; \texttt{USE\_BIOMETRIC} permission
declared; \texttt{MainActivity} extends \texttt{FlutterFragmentActivity}.

\textbf{Web.} PWA; FCM web tokens; PIN fallback for biometric-unavailable
platforms per ADR-0013. WebAuthn foundation per ADR-0014 ships dark.
Responsive breakpoints at 600 / 900 / 1280\,dp.

Single Flutter project builds both targets; \texttt{kIsWeb} branches are
minimal.

\section{Usability Evaluation}

Nielsen heuristic evaluation with two team members. All 10 heuristics
pass with one v1.6 follow-up note (PIN-forgot recovery goes through
account deletion in v1.5; v1.6 plans an email-reset path per ADR-0013).
The single most-valuable evaluation step was the team's read-aloud
session on the Tier 3 curated phrase pool before merge, per ADR-0012
§5.

\printbibliography

\end{document}
